% BOOK-FILE-ID: f-267b502a1839
\subsubsection{桁を取り出す}\label{節：桁を取り出す}
正整数$b\ge2$に対し,非負整数$n$を$b$で割って,
\[
  q=\left\lfloor\frac nb\right\rfloor,
  \qquad r=n-b\left\lfloor\frac nb\right\rfloor
\]
とおくと,\,$n=bq+r,\ 0\le r<b$である.
$b$進法は,数字0から$b-1$を使い,位を$b$倍ずつにする表し方である.
$q$は最後の桁を除いた整数,\,$r$はその桁の数字となる.十進法なら一の位に当たる.
\textbf{床で添字を小さくし,余りで最後の桁を処理する.}

\noindent \bookblocklink{example-gauss-role-8}{problem}{answer}{【例題】}\par\nobreak
\begin{enumerate}[label=(\arabic*),parsep=3pt,beginpenalty=10000,format=\bookitemlink{example-gauss-role-8}{problem}{answer}]
  \item \bookterm{sequence}{数列}を次のように定める.
    \[
      a_0=0,\qquad
      a_n=a_{\left\lfloor\frac n3\right\rfloor}
      +n-3\left\lfloor\frac n3\right\rfloor
      \quad(n\ge1).
    \]
    $a_n$が$n$の3進表示とどのような関係を持つか説明し,\,$a_{14}$を求めよ.
\end{enumerate}

\noindent \bookblocklink{example-gauss-role-8}{hint}{problem}{【方針】}\par\nobreak
\begin{enumerate}[label=(\arabic*),parsep=3pt,beginpenalty=10000,format=\bookitemlink{example-gauss-role-8}{hint}{problem}]
  \item $n=3q+r\ (0\le r<3)$と書くと,\bookterm{recurrence}{漸化式}は$a_n=a_q+r$になります.
    $q$の3進表示に最後の桁$r$を付けると$n$になることを使いましょう.
\end{enumerate}

\noindent \bookblocklink{example-gauss-role-8}{answer}{problem}{【解答】}\par\nobreak
\begin{enumerate}[label=(\arabic*),parsep=3pt,beginpenalty=10000,format=\bookitemlink{example-gauss-role-8}{answer}{problem}]
  \item $a_n$は$n$の3進表示の各桁の数字の和である.
    $14=(112)_3$より,
    \[
      a_{14}=1+1+2=4.
    \]
\end{enumerate}

\noindent \bookblocklink{example-gauss-role-8}{solution}{problem}{【解法】}\par\nobreak
\begin{enumerate}[label=(\arabic*),parsep=3pt,beginpenalty=10000,format=\bookitemlink{example-gauss-role-8}{solution}{problem}]
  \item $n\ge1$では$0\le q=\lfloor\frac n3\rfloor<n$なので,初期値$a_0$から各項を定義できる.
    $r=n-3q$とおけば$r$は0,\,1,\,2のいずれかで,
    \[
      a_n=a_q+r.
    \]
    0の桁の和を0とする.\,$n=3q+r$の表示は,\,$q$の最後に$r$を付けたものなので,桁の和も$r$増える.
    $q<n$と$a_0=0$から,強い帰納法により各項は桁の和に一致する.

    実際,\,$14=3\cdot4+2,\ 4=3\cdot1+1,\ 1=3\cdot0+1$なので,
    \[
      a_{14}=a_4+2=a_1+3=a_0+4=4.
    \]
    値の意味を読むことも,\bookterm{sequence}{数列}を理解する方法である.
\end{enumerate}
